\documentclass{report}
\usepackage{amsmath,amssymb}
\setlength{\parindent}{0mm}
\setlength{\parskip}{1em}
\begin{document}
\begin{center}
\rule{6in}{1pt} \
{\large William W. Dai \\
{\bf On Operator-Splitting Approach for Plasma 3-T Radiation Diffusion in Two and Three Dimensions}}

Los Alamos National Laboratory \\ Mail Stop D413 \\ Los Alamos \\ NM 87545
\\
{\tt dai@lanl.gov}\\
William W. Dai\\
Anthony J. Scannapieco\end{center}

Solving plasma 3-T radiation diffusion equations is a critical step in
multi-physics simulations,
Often these simulations involve mixed cells, and within a mixed cell
there are more than one material, no matter how fine the resolution of a
simulation is. Treatment of mixed cells is important for many
applications. Since materials within a mixed cell could be in a state of
thermally non-equilibrium, and material property of mixture of materials
is unknown, a mixed cell is decomposed to a set of sub-cells through
material interface reconstruction so that each sub-cell contains only one
material. The sub cells thus generated are general polygons in
two-dimensions and polyhedrons in three dimensions. The plasma 3-T
radiation diffusion equations are solved on the mesh with these general
polygons and polyhedrons.


Two typical implicit methods are the backward Euler method and Crank-Nicolson method.
The backward Euler method is first order accurate in time, but numerical
errors in the method undergo quick damping for large time steps.
Therefore the method is very useful for large time steps and steady
states. Although Crank-Nicolson method is second order accurate,
numerical errors do not damp out for large time steps, and significant
numerical errors will be introduced when time steps are large. The
formulation for large time steps is very important for problems involving
dramatically different materials even for time-dependent problems. A
given time step may be considered so large for some material that the
formulation for steady state is more appropriate for the material than
the formulation with the second order of accuracy, but the time step is
so small for other materials that accuracy in time is more important. The
scheme we will present is second order accurate in space and time, works
for any size of time steps, and will give exact steady states when time
steps are very large.

For systems of multi-materials with dramatically different material
properties, the correct treatment for the discontinuity of material
properties is important. A typical approach for this is to use
mathematical approximations, which could introduce numerical errors when
thermal properties of two materials are very different. We use the
governing physics principle to give formula for effective diffusion
coefficient across a material interface for flux calculation on
polyhedral meshes.

In addition to the aspects mentioned above, we will focus on another
issue in numerical simulations for plasma 3-T radiation equations, i.e.,
numerical treatment for interaction between radiation and material that
involves nonlinearity. The 3-T radiation diffusion equations are often
solved through some operator-splitting technique. Several approaches that
are typically used in numerical simulations have serious drawbacks for
some problems.

The numerical scheme to be presented is intended for a code of
multi-physics simulations. It is an extension of previous work to include
mixed cells, unstructured meshes, fully coupling, full nonlinearity, and
operator-splitting. For mixed cells, we consider the interface
reconstruction with any number of materials in both two- and
three-dimensional meshes with adaptive mesh refinement (AMR). In this
talk, we will first describe the procedure for interface reconstruction
for both two- and three-dimensions for AMR meshes with any number of
materials, and then we will present the scheme for the set of equations
for 3-T radiation diffusion in both two- and three-dimensions. We will
demonstrate serious errors of typical operator-splitting approaches for
some problems, and present our operator-splitting techniques.

One operator-splitting approach is to decompose plasma 3-T radiation
diffusion equations into two steps, a set of diffusion equations without
interaction among radiation, electrons and ions, and a set of equations
for the interactions without diffusion. In the first step, the set of
diffusion equations are fully implicitly solved. In the second step, the
coupling among radiation, electrons and ions is also fully implicitly
solved. Through the first step of the operator-splitting approach,
radiation temperature diffuses across the material interface, while
temperatures of electrons and ions may diffuse little. During the second
step of the operator-splitting, almost all the increase of radiation
energy of the heated material transfers to local electrons and ions, but
electrons and ions increase their temperatures very little because of
their large material heat capacity. As a result, three temperatures
barely diffuse through this operator-splitting approach.

A slightly change of this operator-splitting is an approach with three
steps, a half time step of coupling, one time step of diffusion, and a
half time step of coupling. Unfortunately, this operator-splitting
approach will get a similar result.

Another operator-splitting approach for 3-T radiation diffusion equations
is more sophisticated than the one described above, but bears another
serious problem. This approach decompose the set of basic equations into
four steps, each of step is fully implicitly solved. But, due to an
inappropriate treatment for the nonlinear terms, electrons and ions are
much over-heated.

In additional to a fully nonlinear scheme in which material and radiation
are fully coupled, we will present a scheme in which the nonlinearity is
properly treated, and an operator-splitting technique is used, which
generates reasonable solutions for those problems for which the previous
operator-splitting approaches fail.


\end{document}
